\section{The local Lamperti law}\label{sec:local}

\subsection{The limit density}\label{sec:lamperti}

\begin{proposition}[identification of the limit]\label{prop:lambda}
The integral $\Lambda$ of~\eqref{eq:Lambda} is finite and continuous on $(0,1)$, and
$\Lambda=\La$. Moreover, for every $\theta>-1$,
\begin{equation}\label{eq:stieltjes}
 \int_0^1\frac{\La(v)\,dv}{1+\theta v}=\frac{(1+\theta)^{\alpha-1}+b}{(1+\theta)^\alpha+b},
\end{equation}
so $\La$ is a probability density with mean $1/(1+b)$.
\end{proposition}

\begin{proof}
\emph{Finiteness and continuity.} Let $x$ range over a compact subset $K$ of $(0,1]$.
Split $Q_t(x)$ at $z=x/2$. The part $z<x/2$ is at most $(x/2)^{-\alpha}$, and by
\eqref{eq:p-bound} the part $z\ge x/2$ is at most
$C\min\{t^{-1/\alpha},t\}\int_0^{x/2}w^{-\alpha}dw$. So
$Q_t(x)\le C_K(1+\min\{t^{-1/\alpha},t\})$, and the same holds for $Q^b_t$.
By~\eqref{eq:p-bound}, $p_t(x)\le C_K\min\{t^{-1/\alpha},t\}$. The integrand
of~\eqref{eq:Lambda} is therefore at most $C_K\min\{t^{-1/\alpha},t\}$ for $x\in K$,
which is integrable on $(0,\infty)$ because $1/\alpha>1$. The integrand is
continuous in $x$, so $\Lambda$ is finite and continuous by dominated convergence.

\emph{Homogeneity.} For $x,y>0$ put
$\lambda(x,y)=\int_0^\infty[p_{bt}(y)Q_t(x)+p_t(x)Q^b_t(y)]\,dt$, so
$\Lambda(x)=\lambda(x,1-x)$. For $c>0$ the definitions give
$p_t(cz)=c^{-1}p_{tc^{-\alpha}}(z)$, and the substitution $z=cw$ in~\eqref{eq:covering}
gives $Q_t(cx)=c^{-\alpha}Q_{tc^{-\alpha}}(x)$ and $Q^b_t(cy)=c^{-\alpha}Q^b_{tc^{-\alpha}}(y)$.
Substituting $s=tc^{-\alpha}$ in the integral over $t$ then gives
$\lambda(cx,cy)=c^{-1}\lambda(x,y)$. Hence
\[
 \lambda(x,y)=\frac1{x+y}\,\Lambda\Bigl(\frac x{x+y}\Bigr).
\]

\emph{The double Laplace transform.} For $\lambda,\mu>0$,
$\int_0^\infty e^{-\mu y}p_{bt}(y)\,dy=e^{-bt\Gam\mu^\alpha}$, and since $Q_t$ is the
convolution of $p_t$ with $w^{-\alpha}$ on $(0,\infty)$,
\[
 \int_0^\infty e^{-\lambda x}Q_t(x)\,dx=e^{-t\Gam\lambda^\alpha}\,\Gam\lambda^{\alpha-1},\qquad
 \int_0^\infty e^{-\mu y}Q^b_t(y)\,dy=b\,e^{-bt\Gam\mu^\alpha}\,\Gam\mu^{\alpha-1}.
\]
All functions are nonnegative, so by Tonelli's theorem
\[
 \int_0^\infty\!\!\int_0^\infty e^{-\lambda x-\mu y}\lambda(x,y)\,dx\,dy
 =\int_0^\infty\Gam\bigl(\lambda^{\alpha-1}+b\mu^{\alpha-1}\bigr)e^{-t\Gam(\lambda^\alpha+b\mu^\alpha)}\,dt
 =\frac{\lambda^{\alpha-1}+b\mu^{\alpha-1}}{\lambda^\alpha+b\mu^\alpha}.
\]
Now change variables to $x=sv$, $y=s(1-v)$ with $s>0$ and $v\in(0,1)$, so
$dx\,dy=s\,ds\,dv$. By homogeneity the left side is
\[
 \int_0^1\Lambda(v)\int_0^\infty e^{-s(\lambda v+\mu(1-v))}\,ds\,dv
 =\int_0^1\frac{\Lambda(v)\,dv}{\lambda v+\mu(1-v)} .
\]
With $\mu=1$ and $\lambda=1+\theta$ this is~\eqref{eq:stieltjes} with $\Lambda$ in place of
$\La$. At $\theta=0$ it gives $\int_0^1\Lambda=1$.

\emph{Inversion.} For $z\in\C\setminus[0,1]$ let
$G(z)=\int_0^1\Lambda(v)\,dv/(z-v)$, and let
$H(w)=(w^{\alpha-1}+b)/(w^\alpha+b)$ with principal powers on
$\C\setminus(-\infty,0]$. The denominator does not vanish there, since
$|\arg w^\alpha|<\pi\alpha<\pi$. For real $z>1$ put $\theta=-1/z\in(-1,0)$. Then
$1/(z-v)=z^{-1}/(1+\theta v)$, so the transform gives $G(z)=z^{-1}H(1-1/z)$. Since
$1-1/z\in(-\infty,0]$ exactly when $z\in(0,1]$, the right side is analytic on
$\C\setminus[0,1]$, as is $G$. The two agree on $(1,\infty)$, hence on all of
$\C\setminus[0,1]$. Since $\Lambda$ is continuous and integrable,
$-\frac1\pi\mathrm{Im}\,G(x+i\varepsilon)$ tends to $\Lambda(x)$ as $\varepsilon\downarrow0$ for
every $x\in(0,1)$ (Lemma~\ref{lem:poisson}).

We compute the boundary value. Let $x\in(0,1)$ and $w=(1-x)/x$. As
$\varepsilon\downarrow0$, $1-1/(x+i\varepsilon)$ tends to $-w$ from the upper half
plane, so it tends to $we^{i\pi}$ on the principal branch. With $E=e^{i\pi\alpha}$ this
gives $(we^{i\pi})^\alpha=w^\alpha E$ and $(we^{i\pi})^{\alpha-1}=-w^{\alpha-1}E$, so
\[
 G(x+i0)=\frac1x\,\frac{b-w^{\alpha-1}E}{b+w^\alpha E}
 =\frac1x\,\frac{(b-w^{\alpha-1}E)(b+w^\alpha\bar E)}{b^2+2bw^\alpha\cos\pi\alpha+w^{2\alpha}} .
\]
The imaginary part of the numerator is $-b\sin(\pi\alpha)\,w^{\alpha-1}(1+w)$. Hence
\[
 -\frac1\pi\,\mathrm{Im}\,G(x+i0)=\frac{\sin\pi\alpha}\pi\,
 \frac{b\,w^{\alpha-1}(1+w)}{x\,(b^2+2bw^\alpha\cos\pi\alpha+w^{2\alpha})}.
\]
Since $1+w=1/x$, we have $w^{\alpha-1}(1+w)/x=x^{-1-\alpha}(1-x)^{\alpha-1}$, and the
denominator is $x^{-2\alpha}[b^2x^{2\alpha}+2b\cos(\pi\alpha)x^\alpha(1-x)^\alpha+(1-x)^{2\alpha}]$.
So the right side is $\La(x)$, and $\Lambda=\La$. Finally, differentiating
\eqref{eq:stieltjes} at $\theta=0$ gives $-\int_0^1v\La(v)\,dv=-1/(1+b)$.
\end{proof}

The mean $1/(1+b)$ has the expected direction. If $b<1$, the plus runs are heavier
and the plus steps occupy more than half of the time on average.

\subsection{Proof of Theorem~\ref{thm:local}}\label{sec:local-proof}

We first describe the structure of the proof. By the fresh-plus identity, $P_N(k)$ is a sum
over the number $m$ of completed minus runs (first sum) or plus runs (second sum) of a product
of a point probability and a covering probability. Typical configurations have $m$ of order
$1/T_\plus(N)$, since $m$ runs of a law with tail $T$ have total length of order $a_m$, and
$a_m\asymp N$ exactly when $mT_\plus(N)\asymp1$. On this range Gnedenko's theorem and
Lemma~\ref{lem:kernel-limit} turn the sum into a Riemann sum for $\Lambda$. Large $m$ is
negligible because $N$ steps rarely contain many more runs than $1/T(N)$, which is
Lemma~\ref{lem:sums}(iii). Small $m$ is the delicate range. There a few runs have total length
of order $N$, so one run is of order $N$, and the probability that a sum of few runs lands
exactly on a given point is the quantity controlled by the strong renewal theorem. The key
bound of Lemma~\ref{lem:key} removes the covering probability from this range at the cost of a
constant, and Lemma~\ref{lem:cd}(iv) then says that what is left is negligible if and only if
the strong renewal theorem holds. Since every term is nonnegative, the same comparison runs in
the other direction and gives the necessity.

We prove the theorem for a fresh plus start. Fix
$0<x_1<x_2<1$ and let $k$ range over $[x_1N,x_2N]$. Put $j=N-k$, $x=k/N$ and $y=1-x$.
By Proposition~\ref{prop:fresh-identity},
\[
 P_N(k)=\underbrace{\sum_{m\ge1}u^\minus_m(j)\,q^\plus_{m+1}(k)}_{\mathrm{I}}
 +\underbrace{\sum_{m\ge1}u^\plus_m(k)\,q^\minus_m(j)}_{\mathrm{II}}.
\]
Every term is nonnegative. Fix $\delta\in(0,1)$ and split the sums over $m$ into small
$m<\delta/T_\plus(N)$, medium $\delta/T_\plus(N)\le m\le1/(\delta T_\plus(N))$ and large
$m>1/(\delta T_\plus(N))$. We write $\mathrm I_{\rm s}$, $\mathrm I_{\rm m}$, $\mathrm I_{\rm l}$
and so on. Let
\[
 \Lambda_\delta(x)=\int_\delta^{1/\delta}\bigl[p_{bt}(1-x)Q_t(x)+p_t(x)Q^b_t(1-x)\bigr]\,dt .
\]
By Proposition~\ref{prop:lambda}, $\Lambda_\delta$ increases to $\La$ as $\delta\to0$, and
by Dini's theorem the convergence is uniform on $[x_1,x_2]$.

\emph{Medium $m$.} Every medium $m$ is $m_N(t)$ for $t=mT_\plus(N)\in[\delta,1/\delta]$, and
$a_m=Nt^{1/\alpha}(1+o(1))$ uniformly. By Lemma~\ref{lem:cd}(i),
$Nu^\minus_m(j)=p_{bt}(y)+o(1)$ and $Nu^\plus_m(k)=p_t(x)+o(1)$ uniformly. By
Lemma~\ref{lem:kernel-limit}, $q^\plus_{m+1}(k)=T_\plus(N)[Q_t(x)+o(1)]$ and
$q^\minus_m(j)=T_\plus(N)[Q^b_t(y)+o(1)]$ uniformly. Put
$h(t,x)=p_{bt}(1-x)Q_t(x)+p_t(x)Q^b_t(1-x)$. Since $p_t$, $Q_t$ and $Q^b_t$ are bounded on
the relevant compact sets, the products of the error terms are $o(1)$ uniformly, so
\[
 N(\mathrm I_{\rm m}+\mathrm{II}_{\rm m})=\sum_{\delta/T_\plus(N)\le m\le1/(\delta T_\plus(N))}
 T_\plus(N)\bigl[h(mT_\plus(N),x)+o(1)\bigr].
\]
The number of terms is at most $1/(\delta T_\plus(N))$, so the error terms add up to $o(1)$.
The function $h$ is continuous on $[\delta,1/\delta]\times[x_1,x_2]$. Indeed $g$ is
continuous and bounded, so $(t,z)\mapsto p_t(z)$ is continuous on $(0,\infty)^2$, and
$(t,x)\mapsto Q_t(x)$ is continuous by dominated convergence, after the substitution
$z=xw$ in~\eqref{eq:covering}, which gives
$Q_t(x)=x^{1-\alpha}\int_0^1p_t(xw)(1-w)^{-\alpha}dw$. So $h$ is uniformly continuous there,
and the sum is a Riemann sum with mesh $T_\plus(N)$ that converges to $\int_\delta^{1/\delta}h(t,x)\,dt$
uniformly in $x$. Hence
\begin{equation}\label{eq:medium}
 \sup_{x_1N\le k\le x_2N}\bigl|N(\mathrm I_{\rm m}+\mathrm{II}_{\rm m})-\Lambda_\delta(k/N)\bigr|\to0 .
\end{equation}

\emph{Large $m$.} By Lemma~\ref{lem:key} and (R2), $q^\plus_{m+1}(k)\le CT_\plus(k)\le C_{x_1}T_\plus(N)$.
By Lemma~\ref{lem:sums}(iii) for the minus sign, $u^\minus_m(j)\le Ca_m^{-1}e^{-mT_\minus(j+1)/3}$, and
$T_\minus(j+1)\ge cT_\plus(N)$. For $m\ge1/T_\plus(N)$ we have $a_m\ge cN$. Hence
\[
 N\,\mathrm I_{\rm l}\le C\,T_\plus(N)\sum_{m>1/(\delta T_\plus(N))}e^{-cmT_\plus(N)}\le C\,e^{-c/\delta}.
\]
For $\mathrm{II}_{\rm l}$, Lemma~\ref{lem:decay} gives $q^\minus_m(j)\le CT_\plus(N)e^{-cmT_\plus(N)}$,
and $u^\plus_m(k)\le C/a_m\le C/N$, so $N\,\mathrm{II}_{\rm l}\le Ce^{-c/\delta}$ as well.

\emph{Small $m$, sufficiency.} By Lemma~\ref{lem:key}, $q^\plus_{m+1}(k)\le C_{x_1}T_\plus(N)$ and
$q^\minus_m(j)\le CT_\minus(j)\le C_{x_2}T_\plus(N)$. So
\[
 N\,\mathrm I_{\rm s}\le C\,NT_\plus(N)\sum_{m<\delta/T_\plus(N)}u^\minus_m(j),\qquad
 N\,\mathrm{II}_{\rm s}\le C\,NT_\plus(N)\sum_{m<\delta/T_\plus(N)}u^\plus_m(k).
\]
Since $j\ge(1-x_2)N$, (R2) gives $NT_\plus(N)\le CjT_\minus(j)$ and
$\delta/T_\plus(N)\le C\delta/T_\minus(j)$, so $N\,\mathrm I_{\rm s}\le C\,I_\minus(C\delta,j)$
in the notation of~\eqref{eq:cd41}. In the same way
$N\,\mathrm{II}_{\rm s}\le C\,I_\plus(C\delta,k)$. Assume (SRT$_\plus$) and (SRT$_\minus$). Since
$j$ and $k$ are at least $\min(x_1,1-x_2)N$, Lemma~\ref{lem:cd}(iv) gives
\[
 \lim_{\delta\to0}\limsup_{N\to\infty}\sup_{x_1N\le k\le x_2N}N(\mathrm I_{\rm s}+\mathrm{II}_{\rm s})=0 .
\]
Together with~\eqref{eq:medium}, the bound on large $m$ and the uniform convergence of
$\Lambda_\delta$, letting $N\to\infty$ and then $\delta\to0$ gives
$\sup_{x_1N\le k\le x_2N}|NP_N(k)-\La(k/N)|\to0$. This proves (iii)$\Rightarrow$(ii), and
(ii)$\Rightarrow$(i) is trivial.

\emph{Small $m$, necessity.} Assume (i), let $k_N=\lfloor x_0N\rfloor$ and
$j_N=N-k_N$, and take $[x_1,x_2]\ni x_0$. For $m<\delta/T_\plus(N)$ the part $r<k/2$
of~\eqref{eq:q-cover} gives
$q^\plus_{m+1}(k)\ge\Prob(S^\plus_m<k/2)\,T_\plus(k)$, and by Lemma~\ref{lem:sums}(i)
$\Prob(S^\plus_m\ge k/2)\le CmT_\plus(k/2)\le C_{x_1}\delta$. So
$q^\plus_{m+1}(k)\ge\frac12T_\plus(N)$ for $\delta\le\delta_0$. Since every term is
nonnegative,
\[
 \tfrac12NT_\plus(N)\sum_{m<\delta/T_\plus(N)}u^\minus_m(j_N)\le N\,\mathrm I_{\rm s}
 \le NP_N(k_N)-N(\mathrm I_{\rm m}+\mathrm{II}_{\rm m}) .
\]
By (i) and~\eqref{eq:medium} the right side tends to
$\varepsilon(\delta)=\La(x_0)-\Lambda_\delta(x_0)$, and $\varepsilon(\delta)\to0$ as
$\delta\to0$. The map $N\mapsto j_N$ is nondecreasing with increments $0$ or $1$, so it
takes every large integer value $j$. For $j=j_N$, (R2) gives $jT_\minus(j)\le CNT_\plus(N)$
and $\delta/(CT_\minus(j))\le\delta/T_\plus(N)$. Hence
$\limsup_{j\to\infty}I_\minus(\delta/C,j)\le C\varepsilon(\delta)$, and (SRT$_\minus$) follows
from Lemma~\ref{lem:cd}(iv). The same argument with $\mathrm{II}_{\rm s}$, the bound
$q^\minus_m(j)\ge\Prob(S^\minus_{m-1}<j/2)T_\minus(j)\ge\frac12T_\minus(j)$ for small $m$, and the
sequence $k_N$, which also takes every large value, gives (SRT$_\plus$). This proves
(i)$\Rightarrow$(iii).

For a fresh minus start the identity of Proposition~\ref{prop:fresh-identity} has
$q^\minus_{m+1}$ and $q^\plus_m$ in place of $q^\minus_m$ and $q^\plus_{m+1}$. This index
shift changes none of the limits and none of the bounds, and the proof is the same.
\qed

\begin{corollary}\label{cor:local-cases}
Assume (RV$_\alpha$) and (SPAN) and a fresh start of either sign. If $\alpha>1/2$, the
local Lamperti law (ii) of Theorem~\ref{thm:local} holds. If $\alpha\le1/2$, it holds if
and only if the condition of Lemma~\ref{lem:cd}(iii) holds for both run laws, and
(LB$_\pm$) suffices.
\end{corollary}

\begin{proof}
Combine Theorem~\ref{thm:local} with Lemma~\ref{lem:cd}(iii).
\end{proof}

Without the strong renewal theorem, half of the local law survives. The failure can only show
up as an excess.

\begin{corollary}[the lower half of the local law]\label{cor:lower}
Assume (RV$_\alpha$) and (SPAN), and let the start be fresh of either sign. Then for all
$0<x_1<x_2<1$,
\[
 \liminf_{N\to\infty}\ \inf_{x_1N\le k\le x_2N}\bigl(NP_N(k)-\La(k/N)\bigr)\ge0 .
\]
If (SRT$_\plus$) or (SRT$_\minus$) fails, then
$\limsup_NNP_N(\lfloor xN\rfloor)>\La(x)$ for every $x\in(0,1)$.
\end{corollary}

\begin{proof}
The medium and large ranges of the proof of Theorem~\ref{thm:local} use neither (SRT$_\plus$)
nor (SRT$_\minus$). Every term of the fresh-plus identity is nonnegative, so
$NP_N(k)\ge N(\mathrm I_{\rm m}+\mathrm{II}_{\rm m})=\Lambda_\delta(k/N)+o(1)$ uniformly
by~\eqref{eq:medium}, and $\Lambda_\delta\to\La$ uniformly on $[x_1,x_2]$ as $\delta\to0$. This
gives the first claim. If $\limsup_NNP_N(\lfloor xN\rfloor)\le\La(x)$ for some $x$, the first
claim gives convergence at $x$, and Theorem~\ref{thm:local} gives (SRT$_\pm$).
\end{proof}

Theorem~\ref{thm:local} transfers at once to mixed fresh starts, such as the fair start, since
$P_N$ is then a convex combination of $F^\plus_N$ and $F^\minus_N$. It also gives the local law
of the walk itself. The position after $N$ steps is $W_N=2A_N-N$. So for a fresh or mixed
fresh start under (RV$_\alpha$), (SPAN) and (SRT$_\pm$),
\[
 N\,\Prob(W_N=v)=\La\Bigl(\frac{1+v/N}2\Bigr)+o(1)
\]
uniformly over $v\equiv N\pmod 2$ with $v/N$ in each compact subset of $(-1,1)$. The limit
density of $W_N/N$ on $(-1,1)$ is $\frac12\La((1+\xi)/2)$, and the factor $\frac12$ is
compensated by the spacing $2$ of the lattice of $W_N$. For $b=1$ the density
$\frac12L_{\alpha,1}((1+\xi)/2)$ is the scaling function of the ballistic L\'evy walk
in~\cite[Eq.~(61)]{zdk2015}, as one checks by putting $|\xi+1|=2x$ and $|\xi-1|=2(1-x)$ there.

\begin{remark}[the global law]\label{rem:global}
Under (SRT$_\pm$), Theorem~\ref{thm:local} gives
$\Prob(sN\le A_N\le s'N)\to\int_s^{s'}\La$ for $0<s<s'<1$ by Riemann sums, and since $\La$
has total mass $1$ this gives the convergence of $A_N/N$ in law. Proposition~\ref{prop:lambda}
identifies the limit without any input from~\cite{lamperti1958}. This is a derivation of the global law
only under (SRT$_\pm$). The convergence in law is the classical
statement~\cite{lamperti1958}, and we do not use it. The same argument shows that
$\Prob(A_N\le\varepsilon_NN)+\Prob(A_N\ge N-\varepsilon_NN)\to0$ for every
$\varepsilon_N\to0$ under (SRT$_\pm$), since the mass of $P_N$ on $[\delta N,(1-\delta)N]$
tends to $\int_\delta^{1-\delta}\La$. So the end layers of Section~\ref{sec:ends}
dominate pointwise but carry no mass in the limit.
\end{remark}

\begin{remark}[where the bivariate structure went]\label{rem:bivariate}
The count $A_N$ is a functional of the bivariate renewal process
$(S^\plus_m,S^\minus_m)$, and one might expect the local law to need a bivariate strong
renewal theorem. The proof shows that it does not. Lemma~\ref{lem:key} bounds the
covering kernel by $CT(k)$ uniformly in $m$, which decouples the two signs for small
$m$. The small-$m$ part of each sign is then the univariate quantity of
Lemma~\ref{lem:cd}(iv), and a single interior point controls both.
\end{remark}
